\cleardoublepage

\chapter{Przegląd badań porównawczych}
\label{cha:StanBadan}

W~rozdziale zestawiono prace empiryczne porównujące REST z~gRPC lub GraphQL, aby określić, które elementy komunikacji zostały już zbadane i~w~jakim stopniu opublikowane wyniki odnoszą się do~aplikacji działającej w~przeglądarce.

% -----------------------------------------------------------------------
\section{Zakres i~kryteria analizy}
\label{sec:KryteriaPrzegladu}
% -----------------------------------------------------------------------

Do~przeglądu włączono publikacje, w~których podano środowisko wykonania, rodzaj obciążenia oraz co~najmniej jedną miarę wydajności: czas odpowiedzi, rozmiar komunikatu albo wykorzystanie zasobów. Przy ich interpretacji rozpatrywano granicę pomiaru, użyte środowisko, funkcjonalną równoważność operacji oraz sposób dostępu do~danych.

Generator obciążenia rejestruje czas obsługi wywołania API, podczas gdy pomiar realizowany w~przeglądarce obejmuje dodatkowo odebranie odpowiedzi, jej dekodowanie oraz aktualizację interfejsu. Obie wartości odpowiadają odmiennym pytaniom badawczym, dlatego zestawiono je w~sposób jakościowy.

% -----------------------------------------------------------------------
\section{Porównania REST i~gRPC}
\label{sec:PorownaniaRestGrpc}
% -----------------------------------------------------------------------

Rozmiar komunikatu i czas operacji nie pozostają w omawianych badaniach w jednoznacznej zależności. Śliwa i~Pańczyk~\cite{bib:SliwaPanczyk2021} zbudowali trzy aplikacje o~tej samej funkcjonalności na~platformie .NET~5, udostępniające REST, GraphQL i~gRPC. Mniejsze komunikaty przesyłało gRPC, a~korzystne czasy odpowiedzi i~przepustowość w~wielu scenariuszach osiągnął REST. U~Niswara i~współautorów~\cite{bib:Niswar2024} układ ten był odwrotny: w~usługach napisanych w~języku Go, obciążanych narzędziem JMeter, przy odczytach obsługiwanych z~pamięci podręcznej Redis najkrótsze czasy odpowiedzi uzyskało gRPC. Rozmiar komunikatu nie wyznacza więc czasu operacji samodzielnie.

Skoro kierunek relacji zmienia się między badaniami, pozostaje pytanie, czemu właściwie przypisać zmierzoną różnicę. Bolanowski i~współautorzy~\cite{bib:Bolanowski2022} oraz Kamiński i~współautorzy~\cite{bib:Kaminski2023} wiążą wyniki z~rodzajem operacji, środowiskiem uruchomieniowym i~konfiguracją transportu. W~porównaniach tego rodzaju zmieniają się równocześnie format reprezentacji, wersja HTTP, biblioteka serwerowa, a~czasem także ścieżka dostępu do~danych. Mierzony jest wówczas kompletny wariant implementacyjny wraz z~otoczeniem.

Osobną kwestią jest to, gdzie kończy się pomiar. Zarówno Śliwa i~Pańczyk, jak i~Niswar ze~współautorami rejestrowali czas obsługi wywołania po~stronie usługi, przy użyciu generatora obciążenia. Dekodowanie odpowiedzi oraz aktualizacja interfejsu pozostały poza obserwacją, a~w~aplikacji przeglądarkowej właśnie ten odcinek dopełnia czas widziany przez użytkownika.

% -----------------------------------------------------------------------
\section{GraphQL a~zakres zwracanych danych}
\label{sec:GraphQLPrzeglad}
% -----------------------------------------------------------------------

GraphQL umożliwia klientowi dobór pól odpowiedzi, przez co wpływa zarówno na~objętość danych, jak i~na~liczbę żądań potrzebnych do~realizacji operacji. Optymalizacja dotyczy w~tym przypadku zakresu przesyłanych informacji, podczas gdy zastosowanie gRPC zmienia sposób kodowania tej samej treści.

Brito, Mombach i~Valente~\cite{bib:Brito2019} ocenili skutki zastąpienia interfejsów REST rozwiązaniem GraphQL i~wykazali, że~ograniczenie zakresu zwracanych pól istotnie zmniejsza rozmiar odpowiedzi. Vadlamani i~współautorzy~\cite{bib:Vadlamani2021} porównali oba rozwiązania na~publicznym API GitHuba, żadne z~nich nie okazało się szybsze dla wszystkich analizowanych par zapytań. W~badaniu Jina i~współautorów~\cite{bib:Jin2025} porównywane rozwiązania wdrożono w~odmiennej infrastrukturze, wobec czego zmierzona różnica obejmuje również skutki wdrożenia. Warunkiem porównywalności pozostaje zwracanie tego~samego zakresu informacji oraz wykonywanie tej~samej logiki biznesowej.

Wartości procentowych podawanych w~poszczególnych publikacjach nie zestawiano ze~sobą liczbowo. Tabela~\ref{tab:StanBadan} porządkuje te~cechy omówionych prac, które decydują o~możliwości odniesienia ich wyników do~aplikacji przeglądarkowej.

\begin{table}[htbp]
\centering
\small
\caption{Zakres omówionych badań i~ograniczenia ich odniesienia do~eksperymentu przeglądarkowego}
\label{tab:StanBadan}
\renewcommand{\arraystretch}{1.15}
\begin{tabular}{|p{0.25\textwidth}|p{0.31\textwidth}|p{0.32\textwidth}|}
\hline
\textbf{Publikacja} & \textbf{Przedmiot oceny} & \textbf{Ograniczenie odniesienia} \\
\hline
Śliwa, Pańczyk~\cite{bib:SliwaPanczyk2021} & REST, GraphQL i~gRPC; .NET, generator obciążenia & Czas API, bez aktualizacji DOM \\
\hline
Niswar i~in.~\cite{bib:Niswar2024} & REST, GraphQL i~gRPC; Go, Redis, JMeter & Odczyty z~cache, inne środowisko wykonania \\
\hline
Bolanowski i~in.~\cite{bib:Bolanowski2022} & REST i~gRPC w~mikroserwisach & Wynik zależny od~zadania i~konfiguracji \\
\hline
Kamiński i~in.~\cite{bib:Kaminski2023} & REST, gRPC, GraphQL i~WebSocket & Zależność od~środowiska i~bibliotek \\
\hline
Vadlamani i~in.~\cite{bib:Vadlamani2021} & REST i~GraphQL w~API GitHuba & Mała liczba zapytań, brak kontroli implementacji \\
\hline
Brito i~in.~\cite{bib:Brito2019} & REST i~GraphQL; zakres zwracanych pól & Rozmiar odpowiedzi, bez pomiaru po~stronie klienta \\
\hline
Jin i~in.~\cite{bib:Jin2025} & REST i~GraphQL & Odmienna infrastruktura wdrożeniowa porównywanych rozwiązań \\
\hline
\end{tabular}
\end{table}

% -----------------------------------------------------------------------
\section{Ograniczenia porównywalności i~zakres wkładu pracy}
\label{sec:LukaBadawcza}
% -----------------------------------------------------------------------

Natywne gRPC oraz gRPC-Web wymagają odmiennego sposobu realizacji komunikacji, dlatego z~perspektywy aplikacji przeglądarkowej pozostają rozwiązaniami rozłącznymi~\cite{bib:gRPCWeb,bib:ProtocolWeb}. Sam REST dopuszcza z~kolei różne wersje protokołu HTTP, a~jednoczesna zmiana formatu danych, wersji HTTP, obecności TLS i~serwera pośredniczącego prowadzi do~porównania całych ścieżek komunikacji. Mniejszy rozmiar odpowiedzi pomaga wyjaśnić zaobserwowaną różnicę czasu, choć oszacowanie udziału poszczególnych czynników wymaga odrębnego układu eksperymentalnego.
